\begin{explanationbox}
\textbf{\textcolor{CExplanation}{一句话直觉}}

\textbf{找轴是猜测，验轴是证明。} 猜测来自图形特征，证明来自\emph{任一点}与它的对称点都满足同一个曲线条件。

\medskip
\begin{center}
\begin{tikzpicture}[font=\small,>=Stealth,line cap=round,line join=round]
\path[use as bounding box] (-6.5,-3.55) rectangle (6.5,3.55);
\draw[black!20,rounded corners=3pt] (-6.5,-3.55) rectangle (6.5,3.55);
\node[font=\small\bfseries] at (0,3.08) {特殊轴：交换坐标与交换后取负};

\foreach \dx/\sign/\ttl in {-3.25/1/{$y=x$},3.25/-1/{$y=-x$}} {
\begin{scope}[shift={(\dx,0)},x=.7cm,y=.7cm]
\draw[->,black!40] (-3.4,0)--(3.4,0) node[right] {$x$};
\draw[->,black!40] (0,-3.4)--(0,3.4) node[above] {$y$};
\draw[CProof,dashed] (-3,{-3*\sign})--(3,{3*\sign});
\node[fill=white,inner sep=1pt] at (-2.2,{-2.2*\sign}) {\ttl};
\draw[CExplanation,thick,domain=.34:3.15,samples=80] plot (\x,{1/\x});
\draw[CExplanation,thick,domain=-3.15:-.34,samples=80] plot (\x,{1/\x});
\fill[CExplanation] (2,.5) circle(2pt) node[above right] {$(2,\frac12)$};
\fill[CStatement] ({.5*\sign},{2*\sign}) circle(2pt);
\draw[CStatement,<->] (2,.5)--({.5*\sign},{2*\sign});
\ifnum\sign=1
\node[above left] at (.5,2) {$(\frac12,2)$};
\else
\node[below left=3pt] at (-.5,-2) {$(-\frac12,-2)$};
\fi
\end{scope}}

\node[align=center] at (0,-3.05) {两图均为 $xy=1$：$(x,y)\mapsto(y,x)$；$(x,y)\mapsto(-y,-x)$。};
\end{tikzpicture}
\end{center}


蓝色曲线为 $xy=1$。左图交换两坐标，右图交换后同时取负，乘积都不变。注意：坐标轴是渐近线；两条斜线才是对称轴。一般点的检验见证明二。

\medskip
\textbf{\textcolor{CExplanation}{中心改变，轴也跟着平移}}

\begin{center}
\begin{tikzpicture}[font=\small,>=Stealth,line cap=round,line join=round]
\path[use as bounding box] (-6.5,-3.55) rectangle (6.5,3.55);
\draw[black!20,rounded corners=3pt] (-6.5,-3.55) rectangle (6.5,3.55);
\node[font=\small\bfseries] at (0,3.08) {平移后：对称轴必须经过新中心};

\begin{scope}[x=.73cm,y=.73cm]


\draw[dashed,black!55] (0,-3.4)--(0,3.4);
\node[fill=white,anchor=east] at (-.15,2.8) {$x=2$};

\draw[dashed,black!55] (-5,0)--(5,0);
\node[fill=white] at (-4.3,.3) {$y=3$};
\draw[CProof,thick] (-3,-3)--(3,3);
\node[anchor=west] at (3.5,2.8) {$y=x+1$};
\draw[CStatement,thick] (-3,3)--(3,-3) node[right] {$y=-x+5$};
\draw[CExplanation,very thick,domain=.3:4.8,samples=100] plot (\x,{1/\x});
\draw[CExplanation,very thick,domain=-4.8:-.3,samples=100] plot (\x,{1/\x});
\fill (0,0) circle(2pt);
\node[anchor=west] at (1.4,-.55) {$O'(2,3)$};
\end{scope}

\node[align=center] at (0,-3.05) {以 $O'$ 为中心看 $u=x-2,v=y-3$，即恢复 $uv=1$ 的结构。};
\end{tikzpicture}
\end{center}


令 $u=x-2,v=y-3$，方程变成 $uv=1$。把 $v=\pm u$ 换回原坐标，得到 $y=x+1$、$y=-x+5$。先平移看结构，比直接套通式更省事。

\medskip
\textbf{\textcolor{CExplanation}{三层快判流程}}
\begin{itemize}[leftmargin=*, itemsep=0.45em]
  \item \textbf{第一层：识别结构。} 平移后是否是偶函数、$xy=c$、椭圆或双曲线？观察顶点、中心、焦点与渐近线，提出\emph{候选轴}。
  \item \textbf{第二层：选最省事的代换。} $x=h$ 用 $x\to2h-x$；$y=x$ 用交换；$y=-x$ 用交换并取负；其他直线先平移、必要时使用一般公式。
  \item \textbf{第三层：查全体点。} 证明代换后的方程与原曲线有相同解集；若想否定，只需寻找一个真实曲线点，其对称点不在曲线上。
\end{itemize}

\medskip
\textbf{\textcolor{CExplanation}{两个关键区别}}
\begin{itemize}[leftmargin=*, itemsep=0.45em]
  \item \textbf{点对称不等于轴对称：} $P(x,y)\mapsto(2h-x,2k-y)$ 是关于点 $(h,k)$ 的中心对称，不是关于任意一条直线的轴对称。
  \item \textbf{表达式不变不等于唯一标准：} $F(x',y')=F(x,y)$ 是充分条件；严格的核心是\textbf{零点所表示的点集不变}，即 $F=0\Longleftrightarrow F'=0$。必要时要保留定义域及分母不为零的要求。
\end{itemize}

\medskip
\textbf{\textcolor{CExplanation}{顿悟点}}

\textbf{对称检验不是“代进去凑出某个值”，而是“把整条曲线上的任意点送到一个新位置，看看这个新位置是否还满足同一条件”。} 这与只观察几个特殊点有根本区别。
\end{explanationbox}
